\title{LongRoPE: Extending LLM Context Window Beyond 2 Million Tokens}

\begin{abstract}
Having a large context window is an important characteristic for large language models (LLMs). Yet, existing extended context windows typically do not exceed approximately 128k tokens. This limitation arises from factors such as the substantial expenses associated with fine-tuning, the limited availability of sufficiently long training texts, and the problematic values stemming from the introduction of novel token positions.

This work presents LongRoPE, a novel approach that, for the first time, scales the context window of pre-trained LLMs to a remarkable \textbf{2048k} tokens. Remarkably, this extension can be accomplished with as few as 1k fine-tuning steps and training at sequence lengths up to 256k, all while preserving accuracy at the model’s initial short context window. The advancements rest on three core contributions: (i) we uncover and harness two distinctive types of non-uniformities in positional interpolation via a streamlined search procedure, enabling better fine-tuning initialization and facilitating an eightfold window extension even without fine-tuning; (ii) we propose a staged extension process—initially fine-tuning the model for a 256k context, followed by a second round of positional interpolation on this fine-tuned LLM to push the window to 2048k; (iii) we recalibrate LongRoPE at 8k sequence length to recover its short context window capabilities. Comprehensive experiments on both LLaMA2 and Mistral, spanning diverse benchmarks, validate the utility and robustness of our technique. Critically, models upgraded using LongRoPE retain their original structure except for minor tweaks to positional embeddings, thus remaining compatible with established optimizations. The code will be released at \texttt{https://github.com/microsoft/LongRoPE}

\section{Introduction}

Large Language Models (LLMs) have demonstrated impressive capabilities across a variety of applications~\cite{gpt4,llama2}, yet they are constrained by their relatively small context windows, such as LLaMA2's 4096-token restriction~\cite{llama2}. When the input exceeds this limit, the performance of LLMs deteriorates, as the models have not been trained to handle inputs at those extended positions. This constraint creates difficulties in scenarios requiring substantial in-context learning with many examples~\cite{Huang2023FewerIM}, as well as in applications involving LLM agents~\cite{park2023generative,madaan2023selfrefine}.

Recent studies demonstrate that the context window of a pre-trained LLM can be increased to approximately 128k through fine-tuning on lengthier text sequences~\cite{longlora,pi,yarn,zhang2024soaring,liu2023scaling}. Nevertheless, expanding the context window further is impeded by three primary challenges. Firstly, integrating new, untrained positional indices can generate a considerable number of catastrophic values, resulting in out-of-distribution problems that hinder the fine-tuning process from effectively converging~\cite{pi}. This issue becomes even more pronounced when scaling from a 4k to $>$1000k window, as it introduces more than 90\% previously unseen positions. Secondly, effective fine-tuning typically necessitates access to texts of the extended length, but current datasets contain few samples surpassing the 1000k token threshold. Additionally, training on such extensive inputs incurs substantial computational costs, demanding significant time and GPU resources. Thirdly, as the context window is stretched to these extreme lengths, the model's attention mechanism is distributed across many positions, thereby weakening its focus and subsequently diminishing its performance on tasks that involve the original, shorter contexts~\cite{pi}.

To address the initial issue, one commonly adopted method is to interpolate RoPE positional embeddings~\cite{rope,pi}, which rescales novel position indices back to the pretrained domain, as illustrated in Fig.\ref{fig:overview}. Specifically, Position Interpolation (PI)~\cite{pi} conducts linear interpolation of RoPE’s rotary angles, scaling by the proportion of window extension. In contrast, NTK~\cite{ntk,dynamicntk} proposes dimension-wise unequal interpolation and extrapolation when adjusting RoPE embeddings. YaRN~\cite{yarn} subdivides RoPE dimensions according to frequency, assigning extrapolation, NTK, or linear interpolation to three distinct frequency groups. Nonetheless, within the Transformer’s design, positional embeddings encompass \emph{highly non-uniform information entropy}. Current solutions are unable to adequately utilize these nuanced non-uniformities, resulting in some degree of information degradation and, consequently, a restriction on expanding the context window length.

Section~\ref{sec:analysis} identifies two primary empirical insights: \textbf{(1)} For positional interpolation to be effective, it must address two types of non-uniformities—specifically, discrepancies across RoPE dimensions and among token positions. Less interpolation is required for lower RoPE dimensions and for tokens appearing earlier in the sequence; however, the most effective configuration is contingent upon the desired sequence length extension.
  \textbf{(2)} Incorporating these non-uniform characteristics into positional interpolation enables more effective preservation of information within the original RoPE—especially in crucial dimensions and positions—thereby reducing informational loss from interpolation. This leads to superior initialization performance for subsequent fine-tuning. Additionally, it supports sequence extensions of up to 8$\times$ without the need for further fine-tuning.

Building upon these observations, we present \textbf{LongRoPE}, a robust approach designed to broaden the LLM context window to exceed 2 \emph{million} tokens. LongRoPE integrates three principal innovations. First, it leverages the multidimensional variability inherent in positional interpolation to its fullest extent. Specifically, LongRoPE determines optimal rescale factors for RoPE’s rotation angles independently across each RoPE dimension, in relation to token locations. Given that the space for possible rescale factors expands exponentially with increasing extension ratios, LongRoPE adopts an evolutionary search strategy enhanced by two optimization methodologies to significantly improve search performance. An illustration of the resulting rescaled RoPE configuration produced by this search is provided in Fig.~\ref{fig:overview}.

Subsequently, {LongRoPE} employs a resource-efficient, stepwise extension method to expand the context window to 2048k, circumventing the need for direct fine-tuning on ultra-long texts, which are rare and often inaccessible. The process initiates by identifying a suitable 256k context length for the pre-trained LLM and performing fine-tuning specifically at this length. Following this, owing to our non-uniform positional interpolation that supports up to an 8$\times$ extension in settings without fine-tuning, a secondary search is carried out to determine optimal new RoPE rescale parameters using the LLM that has already undergone fine-tuning at the extended length. Through this methodology, the framework successfully enables a 2048k context window for both LLaMA2 and Mistral~\cite{mistral}.

Lastly, to prevent a drop in performance when processing the original (short) context window, {LongRoPE} further tunes the RoPE rescale factors within the extended LLM. Analogous to the scaling procedure applied for extending from 256k to 2048k, we also scale down to 4k and 8k context windows using the 256k fine-tuned LLM, employing our search algorithm to promote reduced positional interpolation. At inference time, if the sequence length falls below 8k, we replace RoPE with the rescale factors determined through our search process.

Comprehensive experiments conducted with multiple LLMs and a range of long-context tasks highlight the efficacy of our approach. Our findings reveal that LongRoPE consistently preserves low perplexity across evaluation lengths spanning from 4k to 2048k tokens, achieves passkey retrieval accuracy exceeding 90\%, and matches the accuracy of existing methods on widely-used benchmarks within the 4096-token context window. Notably, LongRoPE is broadly compatible with any LLMs utilizing RoPE embeddings. We intend to make both our code and the LongRoPE-2048k models publicly available.

\section{Non-uniformity in Positional Interpolation}
\label{sec:analysis}

\subsection{Preliminary}

Transformer models require explicit positional information, often in the form of position embedding, to represent the order of input tokens. Our work focuses on the RoPE~\cite{rope} position embedding, which is widely used in recent  LLMs. For a token at position index $n$, its corresponding RoPE encoding can be simplified as follows:

\begin{equation}
		\fontsize{7.8}{7.8} \selectfont
	\label{eq:rope}
	\left[ cos(n\theta_0), sin(n\theta_0), cos(n\theta_1), \cdot\cdot\cdot, cos(n\theta_{d/2-1}), sin(n\theta_{d/2-1}) \right]   
\end{equation}

where $d$ denotes the embedding dimension, 
 $n\theta_i$ indicates the rotary phase for the token located at position $n$, 
and $\theta_i=\theta^{-2i/d}$ specifies the rotation frequencies. In RoPE, the typical default for $\theta$ is 10000. 

\textbf{\emph{Context window extension ratio $s$} and positional interpolation}. We introduce $s$ as the factor by which the context length is scaled, defined as the ratio of the new context window length $L'$ to the initial length $L$, i.e., $s = \frac{L'}{L}$.

To extend the context window from $L$ to $L'$, current positional interpolation methods suggest downscaling rotation frequencies $\theta_i$ based on extension ratio $s$. Let $\beta=\theta^{2/d}$, and $\lambda$ denote the actual rescale factor related to $s$, we   unify these positional interpolation methods as follows:

\begin{equation}	
	\fontsize{7.5}{7.5} \selectfont
	\label{eq:unified_rope}
	\begin{aligned}
		\left[cos\left(\frac{n}{\lambda(\beta)^0}\right), sin \left(\frac{n}{\lambda(\beta)^0}\right), cos\left(\frac{n}{\lambda(\beta)^1}\right), \cdot\cdot\cdot,  sin \left(\frac{n}{\lambda(\beta)^{d/2-1}}\right) \right]   
	\end{aligned}
\end{equation}

\textbf{Linear positional interpolation (PI)}. PI~\cite{pi} employs a linear scaling of positional indices, constrained within the sequence length observed during pre-training. Given an intended extension ratio $s$, every position’s rotation angle is uniformly scaled down by a factor of $\lambda=s$ for each dimension in RoPE. Nevertheless, this approach condenses positional signals, making it difficult for the model to differentiate between tokens that are near one another in the sequence. As a consequence, PI often exhibits diminished performance when utilized at larger extension ratios.

\textbf{NTK-based interpolation and extrapolation}. Approaches discussed in ~\cite{ntk,dynamicntk} reinterpret RoPE through the lens of information representation and employ Neural Tangent Kernel (NTK) theory~\cite{jacot2018neural,tancik2020fourier} to guide their designs. To address the problem of position overcrowding inherent to PI, these works propose a strategy that redistributes the effects of interpolation across the various dimensions of RoPE. Specifically, they attenuate the scaling for dimensions corresponding to lower frequencies, while amplifying it for dimensions associated with higher (or, more precisely, lower-frequency) components. This methodology enables both positional interpolation and extrapolation using $\lambda=s^{i}$. The enhanced dynamic NTK method~\cite{dynamicntk} further refines this by modulating the extension ratio as a function of the present sequence length for each position. In contrast to PI approaches, which require fine-tuning, NTK-driven techniques are able to expand the context window even without further model tuning. However, these methods are typically constrained to a maximal extension ratio of about 4$\times$.

\textbf{YaRN}~\cite{yarn} offers notable advancements in positional interpolation. Its methodology segments RoPE dimensions into three distinct groups according to their frequency components, assigning separate interpolation techniques to each. Dimensions with high frequencies are subject to extrapolation ($\lambda$=1), whereas those with low frequencies utilize linear interpolation (PI). Intermediate frequency dimensions are assigned to the NTK method. YaRN's strength rests in this frequency-based partitioning of RoPE dimensions, a process that at present is informed by manual, empirically driven experimentation. As a consequence, the approach may not yield optimal results when applied to newly developed LLMs.

\subsection{Study on Non-uniform Positional Interpolation}

Drawing motivation from NTK and YaRN, we observe that their improvements stem largely from incorporating non-linearities, particularly through addressing diverse frequencies within RoPE dimensions to enhance both interpolation and extrapolation. Nevertheless, these existing non-linear approaches predominantly depend on hand-crafted design choices. This prompts two pertinent inquiries: (1) Does present positional interpolation represent the best possible solution? (2) Might there exist non-linear strategies that remain undiscovered?

In order to address these questions, we employ evolution search (refer to Sec.~\ref{sec:method}) to identify improved non-uniform positional interpolation schemes for LLaMA2-7B. The optimization process leverages perplexity as the evaluation criterion and utilizes five randomly selected samples from the PG19~\cite{pg19} validation dataset. Our experimental study yields the following principal observations.

\textbf{Finding 1}: \textit{RoPE dimensions exhibit substantial non-uniformities, which are not effectively handled by current positional interpolation methods.}

We determine the optimal value of $\lambda$ for each RoPE dimension as specified in Eq.~\ref{eq:unified_rope}. Table~\ref{tbl:exp1} presents a comparison of perplexity scores for LLaMA2-7B across various methods, evaluated on the PG19 and Proof-pile~\cite{proof-pile} test datasets without any fine-tuning. Our optimized configuration yields notably improved results, indicating that both traditional linear (PI) and alternative non-uniform approaches (Dynamic-NTK and YaRN) do not achieve optimal performance. Of particular interest, YaRN exhibits lower effectiveness than PI and NTK on the PG19 dataset, primarily because it does not achieve the intended context window length in LLMs that have not undergone fine-tuning. For instance, in the case of an 8k context length, YaRN's perplexity rises sharply after the 7k token mark.

In our investigation, we observed that the rescaling factors $\lambda$ in Eq.~\ref{eq:unified_rope} vary, in contrast to the constant scale $s$ utilized in the computations for PI, the NTK method, and YaRN's group-level scaling. These variable factors yield marked improvements in LLaMA2’s language modeling, as measured by perplexity, across 8k and 16k token context lengths—even in the absence of fine-tuning. This enhancement can be attributed to the modified positional embedding’s ability to largely retain the characteristics of the original RoPE, with a particular emphasis on maintaining the information of key dimensions, thereby alleviating the model’s challenge in differentiating between closely spaced token positions.

\textbf{Finding 2}: \textit{RoPE for the initial tokens in the input sequence should be extrapolated with less interpolation.} 

For the first $\hat{n}$ tokens in the input sequence, we propose that their RoPE should perform minimal interpolation. The rationale is that these initial tokens garner high attention scores, playing a critical role in the functioning of attention layers—a phenomenon previously noted in Streaming LLM~\cite{streamingllm} and LM-Infinite~\cite{han2023lminfinite}. To test this hypothesis, we increase the context window size to 8k and 16k using both PI and NTK methods, leaving the first $\hat{n}$ (where $\hat n$ = 0, 2, ..., 256) tokens unaltered by interpolation. When $\hat n = 0$, the setup reduces to the standard PI and NTK approaches. Table~\ref{tbl:tokenposition} reveals two main findings: \textbf{(1)} excluding the leading tokens from position interpolation contributes to noticeable performance gains, and \textbf{(2)} the best choice of $\hat n$—the count of these excluded tokens—varies with the desired context extension length.

\textbf{Finding 3}: \textit{Non-uniform positional interpolation effectively extends LLM context window in both fine-tuning and non-fine-tuning settings.} 

Our findings indicate that while our optimized non-uniform positional interpolation markedly enhances performance at 8k and 16k context lengths without additional fine-tuning, further extending the context window necessitates fine-tuning. Accordingly, we perform fine-tuning on LLaMA2-7B using the discovered RoPE interpolation scheme, targeting a 64k context length (details are provided in the Appendix). As demonstrated in Table~\ref{tbl:finetune}, our approach surpasses both PI and YaRN significantly, irrespective of whether LLaMA2-7B has been fine-tuned. This improvement can be attributed to our strategic exploitation of non-uniform positional interpolation, which reduces information degradation and offers superior initialization for the fine-tuning process.

\textbf{Summary}. The present work identifies two types of non-uniformities: variations in RoPE dimensions and differences across token positions. By leveraging these non-uniform aspects in positional interpolation, we achieve substantial improvements in large language model (LLM) context length extension.

\section{LongRoPE}

\label{sec:method}
Inspired by these observations, we propose LongRoPE, which initially applies an effective search strategy designed to harness both forms of non-uniformity. This approach is subsequently employed to expand the LLM context window to exceed 2 million tokens.

\subsection{Problem Formulation}

These two forms of non-uniformity expand the solution space considerably and increase the difficulty of the optimization process. To tackle this issue, we recast the multidimensional non-uniform position interpolation optimization challenge as a search problem.

For a LLM targeting a  context window size of $L'$ and lengthy input documents $\mathbf{X}$, where each $\mathbf{x}\in\mathbf{X}$ surpasses $L'$ in token length, we denote the original rotary angle of the $i^{th}$ dimension in RoPE embedding at token position $n$ as  $\frac{n}{\beta^i}$. The optimization problem is then formulated as follows:

\begin{equation}
\begin{gathered}
		\label{eq:problem}

\underset {\mathbf{x}\in\mathbf{X}; \, |\mathbf{x}|\ge L'}{ \operatorname {arg\,min} } \,  \mathcal{L}\, \left( \text{LLM} (\text{RoPE}, \mathbf{X})\right),\;\text{such that}
\\
\scalebox{0.93}{
$\underset {\begin{subarray}{l} i=0,\cdot\cdot,\frac{d}{2}-1;\\ n\in[ 0,|\mathbf{x}|);\end{subarray}}{\text{RoPE}_{(n)}}= {\left[\cdots,\cos\left(\mathbb{I}(\hat{\lambda}_{i}, \hat{n}) \cdot \frac{n}{\beta^i}\right),\;  \sin\left(\mathbb{I}(\hat{\lambda}_{i}, \hat{n}) \cdot \frac{n}{\beta^i}\right),\cdots\right] }$}\\
	\text{with} \;\; \mathbb{I}(\hat{\lambda}_{i},\hat{n})=\begin{cases}
	1&\text{if $n<\hat{n}$}\\
	{\frac{1}{\lambda}_{i}}&\text{if $n\geq\hat{n}$}
\end{cases}
	\end{gathered}
\end{equation}

In this approach, we define a collection of rescale factors, $\mathbb{I}(\hat{\lambda}_{i},\hat{n})$, to accommodate two distinct types of non-uniformity. Here, $\hat{\lambda_i}$ and $\hat{n}$ represent the irregularities associated with RoPE dimensions and with specific token positions, respectively. The factor $\mathbb{I}(\hat{\lambda}_{i},\hat{n})$ serves to adjust the rotation angle in the $i^{th}$ RoPE dimension, utilizing $\hat\lambda_{i}$ as the rescaling parameter and $\hat{n}$ as a positional threshold for the tokens. For the initial token positions (from $1$ up to $\hat{n}-1$), the adjustment factor $\hat\lambda_{i}$ is not applied; thus, the standard RoPE rotary angle $\frac{n}{\beta^i}$ remains unchanged. Once the position $n$ satisfies $n \ge \hat{n}$, however, the rescaling factor becomes active.

For a desired context window length of $L'$, our goal is to determine the most suitable set of rescaling factors ($\mathbb{I}(\hat{\lambda}_{0},\hat{n})$, $\mathbb{I}(\hat{\lambda}_{1},\hat{n})$, ... $\mathbb{I}(\hat{\lambda}_{i},\hat{n})$, ...), corresponding to each RoPE dimension from the $1^{st}$ through the $d^{th}$. This adjustment aims to enable the target $\text{LLM}$, when using the modified $\text{RoPE}$, to attain the lowest possible next-token prediction loss $\mathcal{L}$ (that is, minimize perplexity) across input samples $\mathbf{X}$ of token length $L'$.

\subsection{Searching the Non-uniform Position Interpolation}

To address the issue outlined in Eq.~\ref{eq:problem}, we propose a straightforward and robust approach that identifies the optimal RoPE rescale factors, thereby maximizing the benefits of multidimensional variations in position embedding.

\textbf{Search space}. Our approach constructs an extensive search space that accommodates both identified non-uniformities. The composition of this search space is presented in Table~\ref{tbl:searchspace}. In particular, we introduce the possibility of assigning a unique rescaling factor for each dimension in the RoPE embedding. For the sake of simplifying the overall design, the search is conducted over $\lambda_i$ and $\hat{n}$, rather than directly over $\mathbb{I}(\hat{\lambda}_{i},\hat{n})$, where the relationship $\hat{\lambda}_{i}=1/\lambda_i$ holds. According to Table~\ref{tbl:searchspace}, $\lambda_i$ is permitted to vary from a minimum of 1.0 (representing direct extrapolation) up to a maximum of $s\times1.25$ (indicating an interpolation broader than PI), with increments of 0.01. Here, $s$ denotes the extension ratio for the intended context window.

The parameter $\hat{n}$ determines how many leading token positions preserve their original positional encoding, foregoing any interpolation (meaning they retain the original RoPE embedding). In practice, we consider $\hat{n}$ values chosen from the set \{0, 1, 2, 4, 8, 12, 16, 20, 24, 28, 32, 64, 128, 256\} during our experiments. If $\hat{n}=0$, this implies that the rescale factors found through the search are applied to all token positions.

\textbf{Evolution-based search}. The search space detailed in Table~\ref{tbl:searchspace} contains a vast array of options for positional interpolation, making effective exploration highly complex. For instance, an extension ratio of $s=4\times$ results in $400^{128/2}\times14=4\times10^{167}$ different possible configurations. As the extension ratio increases, the size of the search space grows exponentially. To tackle this complexity, we employ evolution search~\cite{guo2020single}, incorporating two additional optimization strategies to significantly improve the efficiency of the search process. The full search workflow is described in Algorithm~\ref{alg:search}.

\textit{Optimized initial population generation}. Rather than relying solely on random initialization for the population of $P$ rescale factors, we ensure that the initial population explicitly includes the three RoPE rescale factors associated with PI, NTK, and YaRN. The rest of the $P-3$ individuals are produced by randomly mutating these three rescale factors, each mutation occurring with probability $p$.

\textit{Monotonically non-decreasing constraint}. Once the initial population is produced, we evaluate each individual’s LLM perplexity. This involves implementing the designated RoPE rescale factors on the target LLM and assessing the perplexity of the input $\mathbf{X}$. The highest-ranking $k$ individuals are then selected as parents for the evolutionary process. Yet, due to the expansive nature of the search space, unselective mutation and crossover operations may direct the search toward suboptimal solutions, resulting in many redundant perplexity evaluations. This challenge becomes increasingly significant for large $L'$, since each perplexity inference incurs substantial computational cost.

To tackle this issue, we enforce a monotonicity condition such that the sampled RoPE rescaling coefficients are non-decreasing: $\lambda_i \leq \lambda_{i+1}$. Only those RoPE configurations that fulfill this criterion are utilized for evaluating LLM perplexity, which markedly streamlines the search process. More concretely, we stipulate that $\lambda_i$ should monotonically increase as the RoPE dimension grows (with $i$ ranging from 0 to 63). The rationale for imposing this dimension-wise monotonicity comes from NTK theory~\cite{jacot2018neural,tancik2020fourier,ntk}, which implies that lower-dimensional components (corresponding to higher frequencies) warrant less interpolation (smaller $\lambda_i$ values), while higher-dimensional ones (with lower frequencies) permit greater interpolation (larger $\lambda_i$ values).

\begin{algorithm}[t]
	\small
	\caption{The search algorithm}
	\textbf{Input:} target LLM, input samples $\mathbf{X}$,   population size $P$, mutation size $N_1$, crossover size $N_2$, max iterations $\mathcal{T}$, mutate probability $p$\\

	\begin{algorithmic}[1]
		\label{alg:search}
		\STATE Initialize Top-k as $\phi$;	
		\STATE $\text{P}_0$ $\gets$ \textit{Initialize\_population\_with\_optimization}($P$, $\mathbf{X}$, $p$); 
		\FOR{$i$ from $1$ to $\mathcal{T}$}
		\STATE \textit{Compute\_perplexity}(LLM, $\text{P}_{i-1}$, $\mathbf{X}$);
		\STATE Top-k $\gets$ \textit{Update\_Topk}(Top-k, $\text{P}_{i-1}$);
		\STATE $\text{P}_{mutation}$ $\gets$ \textit{Mutation\_with\_mono\_constraint}(Top-k, $N_1$, $p$); 
		\STATE $\text{P}_{crossover}$ $\gets$ \textit{Crossover\_with\_mono\_constraint}(Top-k, $N_2$);
		\STATE $\text{P}_i \gets \text{P}_{mutation} \cup \text{P}_{crossover} \cup \text{Top-k}$;
		\ENDFOR
		\STATE Return the member of Top-k with the minimum perplexity;
	\end{algorithmic}
\end{algorithm}

\textbf{8$\times$ extension without fine-tuning}. Through evolutionary search, our approach selects non-uniform RoPE rescale factors that prioritize retention of crucial dimensions and positions, thereby reducing information loss due to interpolation. As illustrated in Fig.\ref{fig:non-ft}, this strategy successfully increases LLaMA2's context window from 4k to 32k tokens without the need for fine-tuning. By comparison, current alternatives like PI as well as non-uniform NTK and YaRN exhibit significant rises in perplexity after only a 2$\times$ extension.

\subsection{Extending LLM Context Window to 2048K}
\label{sec:extendto2048k}

\textbf{Progressive extension to 2048k}. Here, we propose a strategy to extend the context length of pre-trained LLMs from the conventional 4k up to over 2048k. Our approach, which utilizes non-uniform positional interpolation, is able to realize an 8$\times$ increase in context length without the need for fine-tuning. However, when seeking more substantial extensions—such as scaling by 512$\times$—fine-tuning becomes indispensable. A common approach involves determining optimal RoPE rescaling parameters for the intended 2048k context length, followed by model fine-tuning. This process, nonetheless, is often hampered by the immense computational costs involved. Furthermore, our observations indicate that achieving effective fine-tuning for such extensive context length expansions presents considerable difficulty (refer to Appendix).

Fortunately, LongRoPE demonstrates strong performance on both baseline and further fine-tuned extended LLMs. To this end, we propose an efficient incremental approach that attains the desired 2048k sequence length by performing only 1k fine-tuning steps, utilizing training sequences no longer than 256k tokens.

$\Diamond$ \textit{LongRoPE search for 256k context window in pre-trained LLMs}. Using LLaMA2 as a case study, we perform a search to reach context window lengths of 128k and 256k tokens. At this phase, the scaling factor increases to 32$\times$ and 64$\times$ correspondingly.

$\Diamond$ \textit{Fine-tuning to 256k}. Next, the pre-trained LLM undergoes a fine-tuning process to reach a context window length of 256k. In detail, we initially fine-tune LLaMA2 over 400 steps utilizing the RoPE rescaling factors intended for 128k. Following this, we update the rescaling factors to those for 256k in the obtained checkpoint and perform a further 600 fine-tuning steps. This staged approach is demonstrated to be more effective than commencing fine-tuning directly at 256k.

$\Diamond$ \textit{Scaling up fine-tuned extended LLM to 2048k utilizing LongRoPE search}. In the last step, an additional search procedure is applied to the already fine-tuned LLM at 256k context length. Through this approach, we achieve an exceptionally large context window of 2048k, accomplished without requiring additional fine-tuning steps. This constitutes a final extension factor of $512\times$.

\textbf{Mitigating performance loss in shorter context windows}. When the context window is greatly increased to 2048k tokens, there is a noted decrease in performance within the initial context range. This phenomenon has been previously observed with positional interpolation~\cite{pi}, which compacts the position embeddings for tokens in the shorter context window into a relatively tight region of the embedding space. This compression leads to reduced effectiveness of the language model. Notably, with an extension ratio of 512$\times$, the position representations within the original 4k-token window are highly compressed, further exacerbating the crowding problem.

To address this issue, we conduct an additional evolutionary search on the extended LLM, targeting the adjustment of RoPE rescale factors specifically for shorter context lengths such as 4k and 8k. Since shorter contexts demand less positional interpolation, we lower the upper bound for the $\lambda$ values considered during this search. When performing inference, the LLM automatically updates the relevant RoPE rescale factors to match the current context length.

\section{LongRoPE}

\label{sec:method}
Building on these observations, we propose LongRoPE, a method that begins by deploying an optimized search algorithm specifically designed to harness the two types of non-uniformity. This enables the extension of the LLM context window to lengths exceeding 2 million tokens.

\subsection{Problem Formulation}

These two sources of non-uniformity expand the solution space considerably and add layers of complexity to the optimization process. To tackle this challenge, we reformulate the optimization of multidimensional non-uniform position interpolation as a search problem.

For a LLM targeting a  context window size of $L'$ and lengthy input documents $\mathbf{X}$, where each $\mathbf{x}\in\mathbf{X}$ surpasses $L'$ in token length, we denote the original rotary angle of the $i^{th}$ dimension in RoPE embedding at token position $n$ as  $\frac{n}{\beta^i}$. The optimization problem is then formulated as follows:

\begin{equation}
\begin{gathered}
		\label{eq:problem}

\underset {\mathbf{x}\in\mathbf{X}; \, |\mathbf{x}|\ge L'}{ \operatorname {arg\,min} } \,  \mathcal{L}\, \left( \text{LLM} (\text{RoPE}, \mathbf{X})\right),	\text{where }
\\
\scalebox{0.93}{
$\underset {\begin{subarray}{l} i=0,\cdot\cdot,\frac{d}{2}-1;\\ n\in[ 0,|\mathbf{x}|);\end{subarray}}{\text{RoPE}_(n)}= {\left[\ldots,\cos\left(\mathbb{I}(\hat{\lambda}_{i}, \hat{n})\cdot\frac{n}{\beta^i}\right),  \sin\left(\mathbb{I}(\hat{\lambda}_{i}, \hat{n})\cdot\frac{n}{\beta^i}\right),\ldots\right] }$}
\\
	\text{with } \mathbb{I}(\hat{\lambda}_{i},\hat{n})=\begin{cases}
	1&\text{if }  n<\hat{n}\\
{\frac{1}{\lambda}_{i}}&\text{if } n\geq\hat{n}
\end{cases}
	\end{gathered}
\end{equation}

In this approach, we define a collection of scaling factors, $\mathbb{I}(\hat{\lambda}_{i},\hat{n})$, to accommodate the two types of non-uniformities encountered. Here, $\hat{\lambda_i}$ captures the variation related to RoPE dimensions, while $\hat{n}$ represents the non-uniformity associated with token position. More precisely, $\mathbb{I}(\hat{\lambda}_{i},\hat{n})$ serves to adjust the rotational angle specifically for the $i^{th}$ RoPE dimension, employing $\hat{\lambda}_{i}$ as the scaling factor and $\hat{n}$ as the position threshold. The adjustment through $\hat{\lambda}_{i}$ is only activated for token positions $n \geq \hat{n}$; for token positions preceding the threshold, i.e., those before $\hat{n}$, the rescale factor remains inactive and the standard RoPE rotational angle $\frac{n}{\beta^i}$ is retained.

Assuming a desired context window length of $L'$, our goal is to determine the best set of rescaling factors ($\mathbb{I}(\hat{\lambda}_{0},\hat{n})$, $\mathbb{I}(\hat{\lambda}_{1},\hat{n})$, ...$\mathbb{I}(\hat{\lambda}_{i},\hat{n})$...) for each RoPE dimension, from the $1^{st}$ to the $d^{th}$. The intent is for the modified $\text{LLM}$, after applying the rescaled $\text{RoPE}$, to minimize the next-token prediction loss, $\mathcal{L}$—corresponding to reduced perplexity—when handling input sequences $\mathbf{X}$ consisting of $L'$ tokens.

\subsection{Searching the Non-uniform Position Interpolation}

In order to address the issue presented in Eq.~\ref{eq:problem}, we propose a straightforward and efficient technique that identifies the optimal RoPE rescale factors, thereby allowing for the comprehensive utilization of multidimensional disparities within position embedding.

\textbf{Search space}. We construct an expansive search space that incorporates both forms of non-uniformity. The details of the search space are presented in Table~\ref{tbl:searchspace}. In particular, our approach enables searching for an individual rescale factor tailored to each dimension within the RoPE embedding. To streamline the process of defining the search space, rather than searching over $\mathbb{I}(\hat{\lambda}_{i},\hat{n})$, we instead parameterize the search by $\lambda_i$ and $\hat{n}$, where $\hat{\lambda}_{i}=1/\lambda_i$. As indicated in Table~\ref{tbl:searchspace}, the range of $\lambda_i$ extends from a lower bound of 1.0 (which corresponds to direct extrapolation) up to an upper bound of $s\times1.25$ (allowing for more pronounced interpolation than PI), incrementing in steps of 0.01. Here, $s$ represents the ratio for extending the target context window.

The parameter $\hat{n}$ determines how many of the initial token positions utilize the unaltered RoPE embedding, bypassing position interpolation. In our experiments, we let $\hat{n}$ vary within the set \{0, 1, 2, 4, 8, 12, 16, 20, 24, 28, 32, 64, 128, 256\}. If $\hat{n}=0$, every token position is assigned rescale factors determined by the search process.

\textbf{Evolution-based search}. The search space defined in Table~\ref{tbl:searchspace} encompasses a vast array of positional interpolation strategies, making systematic exploration exceedingly difficult. For instance, when using a $s=4\times$ extension, the number of possible configurations becomes $400^{128/2}\times14$ = 4$\times10^{167}$. As the extension ratio increases, the scale of the search space grows exponentially. To cope with this issue, we leverage evolution search~\cite{guo2020single} and incorporate two optimization strategies that significantly enhance the efficiency of the search process. The complete search algorithm is presented in Algorithm~\ref{alg:search}.

\textit{Optimized initial population generation}. Rather than populating the initial set of $P$ rescale factors through random selection alone, we explicitly insert the three RoPE rescale factors linked to PI, NTK, and YaRN as distinct members of the starting population. The other $P$-3 individuals are produced by applying random mutations to the values of these three rescale factors with probability $p$.

\textit{Monotonically non-decreasing constraint}. Once the initial population has been formed, we evaluate the perplexity of the LLM for every individual. This involves assigning the associated RoPE rescale factors within the target LLM and then measuring perplexity on the input $\mathbf{X}$. The best-performing $k$ individuals are chosen as parents for the next stage of evolution. Due to the immense search space, indiscriminate use of mutation and crossover may sample low-quality candidates, resulting in redundant perplexity evaluations. This inefficiency is exacerbated when $L'$ is large, since each perplexity measurement requires computationally intensive inference.

To tackle this challenge, we enforce a monotonicity constraint in which the rescaled RoPE factors are required to be non-decreasing: $\lambda_i\le \lambda_{i+1}$. We restrict the set of RoPE configurations used in LLM perplexity assessments to only those that satisfy this condition, thereby greatly narrowing the search space. In practical terms, we specify that $\lambda_i$ must increase monotonically as the RoPE dimension index ($i=0,...,63$) increases. This requirement is grounded in the principles of NTK theory~\cite{jacot2018neural,tancik2020fourier,ntk}, which indicate that components at lower dimensions—corresponding to higher frequencies—should involve less interpolation (entailing a smaller $\lambda_i$), while those at higher dimensions with lower frequencies can sustain more interpolation (implying a larger $\lambda_i$).

\begin{algorithm}[t]
	\small
	\caption{The search algorithm}
	\textbf{Input:} target LLM, input samples $\mathbf{X}$, population size $P$, mutation size $N_1$, crossover size $N_2$, maximum iterations $\mathcal{T}$, mutation probability $p$\\

	\begin{algorithmic}[1]
		\label{alg:search}
		\STATE Top-k $\leftarrow \phi$;
		\STATE $\text{P}_0 \leftarrow$ \textit{Initialize\_population\_with\_optimization}($P$, $\mathbf{X}$, $p$);
		\FOR{$i = 1$ to $\mathcal{T}$}
		\STATE \textit{Compute\_perplexity}(LLM, $\text{P}_{i-1}$, $\mathbf{X}$);
		\STATE Top-k $\leftarrow$ \textit{Update\_Topk}(Top-k, $\text{P}_{i-1}$);
		\STATE $\text{P}_{mutation} \leftarrow$ \textit{Mutation\_with\_mono\_constraint}(Top-k, $N_1$, $p$);
		\STATE $\text{P}_{crossover} \leftarrow$ \textit{Crossover\_with\_mono\_constraint}(Top-k, $N_2$);
		\STATE $\text{P}_i = \text{P}_{mutation} \cup \text{P}_{crossover} \cup \text{Top-k}$;
		\ENDFOR
		\STATE Output the member of Top-k with the minimum perplexity;
	\end{algorithmic}
\end{algorithm}

\textbf{8$\times$ extension without fine-tuning}. Through evolutionary search, our approach efficiently discovers optimal non-uniform RoPE rescale factors, maintaining crucial dimensions and positions in order to reduce information degradation due to interpolation. As illustrated in Fig.\ref{fig:non-ft}, this technique enables LLaMA2's context window to be expanded from 4k to 32k tokens without the need for fine-tuning. By comparison, other approaches like PI, as well as non-uniform NTK and YaRN, result in a significant increase in perplexity following a 2$\times$ extension.

\subsection{Extending LLM Context Window to 2048K}
\label{sec:extendto2048k}

\textbf{Progressive extension to 2048k}. Here, we present our approach for expanding the context window of pre-trained LLMs from the standard 4k up to more than 2048k tokens. Our experiments show that leveraging non-uniform positional interpolation permits an $8\times$ increase in context length without necessitating any fine-tuning. However, when aiming for substantially greater extensions (such as 512$\times$), fine-tuning becomes indispensable. One possible strategy is to identify optimal RoPE scaling parameters suitable for the desired 2048k context size, followed by fine-tuning the model. Nonetheless, this procedure encounters obstacles, notably the immense computational resources required. In addition, as reflected in our findings, successfully fine-tuning LLMs for very large extension ratios remains difficult (refer to Appendix).

Fortunately, LongRoPE performs well with both the base model and the fine-tuned variant extended to longer contexts. Consequently, we propose an efficient, incremental approach that enables reaching the desired 2048k target using only 1k fine-tuning iterations, all conducted at a training length capped at 256k.

$\Diamond$ \textit{Scaling pre-trained LLMs to 256k tokens via LongRoPE search}. Using LLaMA2 as a case study, we experiment with expanding the context window to either 128k or 256k tokens. At these lengths, the context window is enlarged by factors of 32$\times$ and 64$\times$, respectively.

$\Diamond$ \textit{Fine-tuning to 256k}. In this stage, we adapt the pre-trained LLM to handle a 256k context window through fine-tuning. The process begins with 400 fine-tuning steps on LLaMA2 utilizing RoPE rescaled factors calibrated for 128k. Subsequently, these RoPE rescaled factors are updated to 256k on the resulting checkpoint, after which we perform an extra 600 steps of fine-tuning. This staged strategy is shown to be more efficient compared to initiating fine-tuning at 256k from the outset.

$\Diamond$ \textit{Expanding fine-tuned extended LLM to 2048k through LongRoPE search}. In the last step, we carry out an additional search process on the 256k-length LLM that has already been fine-tuned. This approach enables the model to achieve a substantially increased context window size of 2048k, requiring no extra rounds of fine-tuning. As a result, the total extension ratio reaches $512\times$.

\textbf{Recovery of performance in shorter context windows}. Upon expanding the context window to an extensive 2048k length, a decline in performance within the initial context window is observed. This phenomenon has been previously documented for positional interpolation~\cite{pi}, which compresses the position embeddings for the original context window into a limited section of a larger space, thereby impairing the effectiveness of the language model. In particular, when the extension ratio reaches 512$\times$, the positional representations of tokens within the standard 4k context window become highly concentrated.

To address this issue, we conduct an additional evolutionary search on the expanded LLM, tuning the RoPE rescale factors specifically for shorter context lengths (such as 4k and 8k tokens). For these reduced lengths, the upper bound for $\lambda$ during the search is lowered, as less positional interpolation is necessary. At inference time, the LLM adaptively modifies the relevant RoPE rescale factors based on the input context.

 \section{Experiments}

\label{sec:exp}
\subsection{Setup}
\textbf{Evaluation Tasks and Models}. LongRoPE is implemented on both LLaMA2-7B and Mistral-7B architectures. The assessment covers three main areas: (1) measuring perplexity scores of large language models with extended contexts when processing lengthy documents; (2) the Passkey retrieval task, designed to evaluate how effectively a model can extract a straightforward passkey from a large quantity of distractor text; and (3) established large language model benchmarks conducted within a restricted 4096-token context window.

\textbf{Fine-tuning}. For LLaMA2, fine-tuning is performed with a learning rate set to 2e-5 employing linear decay, and the training uses a global batch size of 32. The initial phase comprises 400 steps on the Redpajama~\cite{redpajama} dataset, which is divided into 128k-length segments demarcated by BOS and EOS tokens. Following this stage, training is resumed from the resulting checkpoint for an additional 600 steps to reach a 256k context window. The fine-tuning for the 128k context window leverages 8 A100 GPUs with distributed training as described in~\cite{cube}, while scaling up to 16 A100 GPUs is necessary for the 256k context. For Mistral, training is conducted with a fixed learning rate of 1e-6 and a global batch size of 64. The protocol for both 128k and 256k context models aligns with the YaRN~\cite{yarn} configuration, involving 400 steps on the Together Computer's Long-Data Collections~\cite{mistraldata} using sequences of length 16k. All trainings are performed with 4 A100 GPUs.

\textbf{Search}. For target window sizes up to 256k, we adopt the following configuration: $P$=64, $N_1$=$N_2$=16, $p$=0.3, $\mathcal{T}$=40, and select the top-32 candidates for mutation and crossover during each iteration. Perplexity scores are computed using 5 randomly chosen validation samples from PG19, each meeting or exceeding the required context length. For target windows above 512k, we reduce the sizes for population, mutation, and crossover parameters by half. In these cases, perplexity is evaluated using 3 randomly drawn samples from the Pile-Books3~\cite{pile} validation dataset.

\textbf{Baselines}. To achieve a 2048k context length, we conducted fine-tuning of models at the 128k and 256k context sizes. This process resulted in LongRoPE-2048k (ft=128k) for LLaMA2 and LongRoPE-2048k (ft=256k) for Mistral. Our evaluation contrasts these four configurations with cutting-edge context window extension baselines, focusing on open-source LLMs that underwent fine-tuning subsequent to positional interpolation via PI, NTK, and YaRN. Among the models included for comparison are Together-32k~\cite{together}, Code LLaMA~\cite{codellama}, LongLoRA-full-FT-100k~\cite{longlora}, as well as YaRN-LLaMA and YaRN-Mistral~\cite{yarn}.

\subsection{Main Results}

\label{sec:mainresult}
\textbf{Long sequence language modeling up to 256k tokens}. Our initial analysis focuses on assessing performance relative to leading extended LLMs over sequences of up to 256k tokens. To highlight our approach's versatility, we draw on two distinct datasets: the test portions of Proof-pile~\cite{pg19} and PG19~\cite{pile}. Perplexity scores are calculated across different context window sizes, applying a sliding window of length 256. For PG19, the evaluation covers all 100 documents in the test split. In the case of Proof-pile, consistent with the methodology in YaRN~\cite{yarn}, we randomly sample 10 instances, each exceeding a minimum length of 128k tokens.

Table~\ref{tbl:proof-pile} and Table~\ref{tbl:pg19} present a comparison of perplexity scores for LLaMA2 and Mistral models augmented using several interpolation strategies on the Proof-pile and PG19 datasets, respectively. Two major findings are apparent: \textbf{(1)} across increasing evaluation sequence lengths from 4k up to 256k, the extended models consistently exhibit reduced perplexity, demonstrating improved utilization of long-range context. \textbf{(2)} Remarkably, despite operating with a context window that is 16$\times$ larger—a scenario that typically hampers short-range performance—our LongRoPE-2048k models maintain an advantage, achieving lower perplexity than leading baselines up to the 256k length mark.

\textbf{Long sequence language modeling beyond 2000k}. To assess performance on exceptionally lengthy documents, we utilize the Books3~\cite{pile} dataset. For efficient evaluation, we randomly choose 20 books with lengths greater than 2048k, applying a sliding window of 256k tokens.

As indicated in Table~\ref{tbl:books3}, LongRoPE effectively expands the context window of LLaMA2-7B and Mistral-7B to 2048k tokens, while matching or surpassing the perplexity scores of baseline models for input sequences within the 8k to 128k range. We also notice substantial differences in performance at the 2048k scale between LLaMA2 and Mistral. At shorter context lengths, Mistral achieves lower perplexity than the baselines, yet its perplexity rises above 7 for contexts exceeding 256k. LLaMA2’s results largely follow anticipated trends: perplexity generally diminishes with larger context sizes, although there are modest increases at the 1024k and 2048k marks. Notably, for LLaMA2, LongRoPE-2048k shows improved results when fine-tuned at a 256k sequence length compared to 128k, attributable to a reduced secondary extension ratio (8$\times$ instead of 16$\times$). Conversely, Mistral attains higher performance when fine-tuned with a 128k context window. This discrepancy is mainly due to the training regime: for both 128k and 256k fine-tuning settings in Mistral, we employ a 16k training sequence length following YaRN's methodology, which impacts Mistral's ability to further increase the context window after fine-tuning.

\textbf{Passkey retrieval}. Next, we investigate how well models handle varying context window sizes during generation. To this end, we adopt the synthetic passkey retrieval benchmark introduced by~\cite{passkey}. In this benchmark, models are required to extract a randomly embedded five-digit passkey from an extended document. The full prompt template can be found in the appendix. For evaluation, we conduct 10 runs of the retrieval task, randomly positioning the passkey at any location uniformly distributed over the total context length for each trial.

Figure~\ref{fig:passkey} presents a comparison of retrieval accuracy against baseline models. The accuracy of previous approaches sharply declines to zero beyond the 128k token mark. Conversely, even under the demanding scenario of extracting a passkey from millions of tokens, our LongRoPE-LLaMA2-2048k (ft=256k) consistently achieves high retrieval performance ($\ge$90\%) across the entire range from 4k to 2048k tokens. Similarly, LongRoPE-Mistral-2048k (ft=128k) maintains perfect accuracy up to 1800k tokens, before dropping to 60\% at 2048k—a trend that matches the increase in perplexity observed at 2048k in Table~\ref{tbl:books3}.

\textbf{Standard benchmarks within the original context window}. To assess LongRoPE-2048k models under their initial context window, we conduct evaluations employing the Hugging Face Open LLM Leaderboard~\cite{huggingfaceboard}, considering both zero-shot and few-shot scenarios. Specifically, we adopt the following setups: ARC-Challenge~\cite{allenai:arc} with 25 shots, HellaSwag~\cite{zellers2019hellaswag} with 10 shots, MMLU~\cite{mmlu} with 5 shots, and TruthfulQA~\cite{lin2021truthfulqa} in a 0-shot setting.

According to Table~\ref{tbl:commonsense}, our models perform similarly to prior systems on the established benchmark, which was created for shorter context windows, and even surpass the original Mistral by +0.5\% on the TruthfulQA dataset. The LongRoPE-LLaMA2-2048k variant, trained with a 256k context window, displays a marginally higher reduction in performance; however, for the majority of tasks, its results stay within an acceptable margin.

\subsection{Ablation Results}

\textbf{Effectiveness of the second positional interpolation}. Within our progressive extension framework, we apply our search-based method to perform a secondary non-uniform positional interpolation on extended LLMs that have already been fine-tuned. To assess its impact, we execute a series of evaluations using the fine-tuned LLaMA2-256k model. This model is further extended to 512k, 1024k, and 2048k sequences employing both PI and YaRN. As presented in Table~\ref{tbl:secondsearch}, our approach with non-uniform positional interpolation maintains stable perplexity values. In comparison, the perplexity scores for both PI and YaRN rise significantly as the extension ratio increases.

\textbf{Effectiveness of recovery at shorter context lengths.} To address the drop in performance seen at reduced context lengths, we recalibrate the RoPE factors in LongRoPE-2048k utilizing our search procedure. In this process, we limit the maximum scale factors explored by the search to promote reduced interpolation for shorter sequences of 4k and 8k tokens. Table~\ref{tbl:shortperformance} presents a perplexity comparison for LongRoPE-LLaMA2-2048k on Proof-pile at these context lengths, as well as the average benchmark accuracy for LLM tasks. The data clearly indicate notable gains in performance for shorter contexts.

\textbf{Analysis on the two forms of non-uniformities}. To further examine the influence of each non-uniformity, we conduct an ablation study to isolate their contributions to model performance. Two experimental configurations are used: (i) scaling LLaMA2-7B to shorter contexts (16k and 32k tokens) by employing various strategies—namely, PI, RoPE dimension search alone, and a combined search involving both forms of non-uniformity; (ii) expanding our 256k context-length fine-tuned LLaMA2 up to 2048k, applying analogous procedures. Perplexity metrics are gathered prior to any additional fine-tuning. As presented in Table~\ref{tbl:searchdimension}, introducing non-uniformity in the RoPE dimension yields a marked perplexity reduction when compared with the standard linear interpolation of PI. Adjusting for non-uniformity in token position leads to a clear improvement at context lengths of 16k and 32k, while its effect diminishes for the 2048k scenario—potentially due to the exceedingly long sequence. Simply retaining the initial segment of tokens without interpolation proves ineffective in this regime, and we identify this observation as a direction for future exploration.

\section{Related Works}

Besides techniques relying on position interpolation, this section reviews related studies employing alternative strategies.

\textbf{Retrieval-based} methods leverage external memory systems that store extensive historical context, and incorporate retrieval mechanisms to access relevant documents during inference~\cite{longllama,wang2023augmenting,borgeaud2022improving}. Implementing such methods often requires significant alterations to the underlying LLM architectures. By contrast, our proposed approach introduces only minimal adjustments to the positional embeddings, making it substantially more efficient. Furthermore, our method can accommodate a broader range of long-context applications beyond retrieval, including tasks like long document summarization and few-shot learning.

\textbf{Attention-based context window extensions}. In addition to interpolating positional embeddings, other studies have increased the input context length within the constraints of the original LLM context window by altering attention mechanisms~\cite{han2023lminfinite, streamingllm,parallelwindow}. Central to these approaches is the use of innovative attention masking techniques that address the problem of attention explosion from introducing new positions. Notably, these methods are complementary to positional interpolation techniques.

\textbf{Fine-tuning based approaches} investigate strategies for adapting pre-trained LLMs to support longer contexts by altering position embeddings and subsequently fine-tuning on extended sequences. Approaches such as Code LLaMA~\cite{codellama}, LLaMA2 Long~\cite{xiong2023effective}, and ScaledRoPE~\cite{liu2023scaling} utilize a significantly increased base parameter for RoPE, followed by fine-tuning on the desired sequence length. The technique presented here allows adaptation to a wide range of target lengths, including those surpassing 2M tokens. Furthermore, because scaling context lengths beyond 128k typically results in intensive GPU usage, recent solutions like LongLoRA~\cite{longlora} and PoSE~\cite{zhu2023pose} have been introduced to reduce resource consumption. Our approach remains complementary to these efficient fine-tuning methods.

\section{Conclusion}

In this paper, we introduce LongRoPE, a technique capable of extending the context length of large language models (LLMs) to an unprecedented 2048k tokens, all while preserving their performance on standard, shorter context spans. Our approach takes advantage of two specific types of non-uniformities in RoPE positional embeddings, identified through a highly efficient evolutionary search process. This methodology yields two main advantages: it offers superior initialization for downstream fine-tuning, and it allows for up to an 8$\times$ expansion of the context window even without further fine-tuning. Building upon this foundation, we propose a stepwise extension method whereby fine-tuned LLMs with a 256k token window can be further expanded to a 2048k context window without requiring additional fine-tuning. Our extensive experimental results demonstrate the effectiveness of LongRoPE. We anticipate that our LongRoPE-2048k models will unlock a wide range of novel long-context applications and foster further developments in this area.

\section*{Broader Impacts} The aim of this paper is to contribute to ongoing progress in the discipline of Machine Learning. Although our research may have various implications for society, we do not identify any particular impacts that warrant special attention in this context.

	\balance

\end{document}